% ════════════════════════════════════════════
% 具体分析（5.3.1）
% ════════════════════════════════════════════
\subsubsection{建模思路}

本部分属于回归建模与预测解释，采用多元线性回归、HC3 稳健标准误、共线性诊断与交叉验证相结合的方法。与单纯追求拟合优度不同，本部分首先要保证系数可解释，其次才比较预测精度。

首先以 log1p 曝光量为因变量，构建包含粉丝规模、帖子类型、内容主题、月份、星期、时段与付费标记的全模型，并用 HC3 协方差矩阵处理异方差。接着计算方差膨胀因子、广义方差膨胀因子与标准化设计矩阵条件数，定位共线的具体来源，并结合前文的月内方差证据判断哪些参数不可识别。然后根据诊断结果重设模型，剔除粉丝规模、保留月份固定效应，对最终模型做残差正态性、异方差、函数形式与强影响点检验。最后用十折交叉验证比较五种设定的样本外表现，并以 Shapley 分解量化各组变量对解释方差的贡献。具体分析流程如图\ref{fig:analysis3_flow}所示。

\begin{figure}[ht]
    \begin{center}
        \begin{tikzpicture}[x=1cm, y=0.7cm]
            \node[box] (n11) at (0, 0) {构建设计矩阵};
            \node[box] (n12) at (5, 0) {全模型估计};
            \node[box] (n13) at (10, 0) {共线诊断};
            \node[box] (n22) at (5, -3) {模型重设};
            \node[box] (n21) at (0, -3) {残差检验};
            \node[box] (n23) at (10, -3) {整体显著性};
            \node[box] (n31) at (0, -6) {交叉验证};
            \node[box] (n32) at (5, -6) {方差分解};
            \node[box] (n33) at (10, -6) {输出最终模型};
            \draw[arrow] (n11) -- (n12);
            \draw[arrow] (n12) -- (n13);
            \draw[arrow] (n23) -- (n22);
            \draw[arrow] (n22) -- (n21);
            \draw[arrow] (n13) -- ++(0,-1.0) -| (n23);
            \draw[arrow] (n21) -- (n31);
            \draw[arrow] (n31) -- (n32);
            \draw[arrow] (n32) -- (n33);
        \end{tikzpicture}
        \caption{回归建模流程图}
        \label{fig:analysis3_flow}
    \end{center}
\end{figure}
